\infographicchapter{Photoacoustic Imaging (PAI)}{pai.png}
\label{ch:infographic-pai}

\section{Overview}

Photoacoustic imaging (PAI)---also known as optoacoustic imaging---is an emerging hybrid medical scan based on the photoacoustic effect that turns harmless pulses of light into gentle sound waves. When a brief flash of laser light hits natural light-absorbing molecules inside the body, those molecules warm slightly and expand just enough to send out tiny ultrasound waves. Because sound travels through soft tissue much more cleanly than light does, PAI overcomes the main limit of regular optical cameras: it keeps rich light-based contrast while reaching deeper with ultrasound-like sharpness. Depending on the scanner design, PAI can zoom in on microscopic surface capillaries (photoacoustic microscopy, or PAM) or scan several centimeters deep using handheld probes or tomography cups (photoacoustic computed tomography, or PACT).

Without injecting any contrast dye, PAI naturally maps oxygen-rich and oxygen-poor blood (showing tiny blood vessels and oxygen levels), skin pigment (melanin), fat, and collagen. When extra targeting is helpful, it can also detect safe, FDA-approved medical dyes such as indocyanine green (ICG) and methylene blue, as well as experimental nanoparticles being studied in labs.

In everyday medicine, PAI is still mostly a research technology under active clinical testing. One major milestone is U.S. FDA approval of an adjunctive optoacoustic-ultrasound breast scanner that helps doctors tell benign breast lumps from cancerous ones alongside standard ultrasound. Researchers are also testing PAI in clinical trials for melanoma lymph node checks, surgical tumor borders, inflammatory arthritis, Crohn's disease, neck-artery plaques, and poor circulation. Wider hospital use awaits standardized calibration, better depth correction, and larger patient studies.

\section{Underlying Physical Principles}
\begin{itemize}
  \item \textbf{Brief Pulses of Safe Light:} A pulsed visible or near-infrared laser shines billionth-of-a-second flashes of non-ionizing light into tissue, kept strictly within ANSI and IEC safety limits.
  \item \textbf{Light Creates Gentle Sound Waves:} Natural absorbers---such as blood hemoglobin, melanin, fat, and collagen---or approved dyes absorb the light flash, expand slightly, and launch broadband ultrasound waves.
  \item \textbf{Works from Microscopic to Centimeter Scales:} Because sound scatters far less than light in soft tissue, PAI scales smoothly from microscopic capillary imaging near the surface (about 1~millimeter deep) to sub-millimeter tomography down to 2 to 4~centimeters deep.
  \item \textbf{Shared Ultrasound Sensors:} Sensitive ultrasound probes listen for the returning sound waves and can easily share the same handheld transducer used for standard grayscale B-mode ultrasound.
  \item \textbf{Multi-Color Oxygen Mapping (MSOT):} Because blood absorbs light much more strongly than surrounding soft tissue, comparing multiple light wavelengths---an approach often called multispectral optoacoustic tomography (MSOT)---lets software measure total hemoglobin levels and separate oxygen-rich from oxygen-poor blood to map blood oxygen saturation and tiny-vessel density.
\end{itemize}

\section{Imaging Mechanism}
\begin{enumerate}
  \item \textbf{Eye Safety and Probe Coupling:} Both you and the staff wear wavelength-matched laser safety goggles. The clinician applies clear, bubble-free ultrasound gel (or a soft water membrane) and rests the handheld probe or imaging cup gently on the skin.
  \item \textbf{Paired Laser Flashes and Ultrasound:} The pulsed laser fires rapid flashes of visible or near-infrared light into the area, alternating smoothly with standard ultrasound pulses. Microscopy systems scan a focused spot, while handheld and tomography units light up a wider patch at once.
  \item \textbf{Picking Up the Sound Waves:} Ultrasound sensor elements catch the returning sound waves, boost the faint signals, and convert them into digital data.
  \item \textbf{Building and Unmixing the Image:} Reconstruction software turns the sound data into absorption maps and separates out oxygen-rich blood, oxygen-poor blood, or injected dyes while adjusting for depth.
  \item \textbf{Live Dual-Screen Overlay:} Color maps of blood vessels, low-oxygen zones, or dye appear directly on top of the live grayscale ultrasound image for immediate bedside review.
\end{enumerate}

\section{Key Technical Parameters}
\noindent
\begin{tabularx}{\textwidth}{@{}l X@{}}
\toprule
\textbf{Parameter} & \textbf{Typical Value / Description} \\
\midrule
Physical Principle & Pulsed non-ionizing light absorbed by blood and tissue creates gentle ultrasound waves that sensors reconstruct into vascular and oxygen maps \\
Energy / Signal Source & Pulsed tunable visible or near-infrared laser paired with ultrasound transducer arrays (microscopy, tomography, or handheld units) \\
Typical Frequency / Energy & 532--1,064~nm light pulses (within ANSI/IEC safety limits); 2--50~MHz detected sound waves \\
Spatial Resolution & 5--50~micrometers near the surface (PAM) to 0.1--0.5~mm at 2--4~cm clinical depth (PACT and handheld PAI) \\
Typical Scan Duration & 15--30~minutes for a targeted handheld clinical session \\
Contrast Agent & None needed for natural blood, oxygen, melanin, fat, and collagen; optional approved dyes (ICG, methylene blue) or research nanoparticles \\
Ionizing Radiation & None (uses non-ionizing light pulses and sound waves) \\
\bottomrule
\end{tabularx}

\section{Clinical Usages}

\subsection{Primary Usages}
\begin{itemize}
  \item \textbf{Helping Tell Benign from Malignant Breast Lumps (FDA-Cleared Adjunct):} An FDA-approved handheld optoacoustic-ultrasound breast system pairs blood-vessel and oxygen maps with regular ultrasound without compressing the breast. Clinical studies involving over 2,000~patients showed that its vessel, hemoglobin, and border-zone scores improve radiologists' accuracy when evaluating solid breast lumps, while larger dedicated 3D whole-breast bowl scanners are also being studied.
  \item \textbf{Skin Cancer and Melanoma Staging (Research Target):} Pilot clinical studies test PAI to measure melanoma depth before surgery, check nearby lymph nodes for melanoma pigment without cutting, and check tumor edges during breast or head-and-neck surgery.
  \item \textbf{Spotting Early Joint Inflammation (Clinical Trials):} Handheld probes are being tested in rheumatoid and psoriatic arthritis to detect early joint-lining swelling and low oxygen before standard Doppler ultrasound picks it up.
  \item \textbf{Checking Small Blood Vessels (Research Target):} PAI is being studied for Raynaud's phenomenon, diabetic microvascular disease, slow-healing wounds, and vulnerable neck-artery plaques.
\end{itemize}

\subsection{Secondary Usages}
\begin{itemize}
  \item \textbf{Tracking Bowel Inflammation (Investigational):} External and endoscopic PAI prototypes are being tested to measure bowel-wall oxygen and scarring in Crohn's disease and inspect digestive-tract linings.
  \item \textbf{Guiding Needles and Catheters (Research):} Researchers are exploring PAI to track biopsy needles, radiation seeds, and catheters while clearly highlighting nearby blood vessels and nerves.
  \item \textbf{Early Thyroid, Prostate, and Ovarian Studies:} Pilot trials pair neck or internal ultrasound probes with PAI to look at tumor blood supply and oxygen in deeper organs.
  \item \textbf{Measuring Muscle Collagen Without Sedation (Pilot Trials):} Exploratory studies use multispectral PAI to track muscle collagen buildup in Duchenne muscular dystrophy in under 7~minutes across multiple muscles, offering a fast, dye-free alternative when lengthy MRI scans might require sedation.
  \item \textbf{Preclinical Lab Research:} Photoacoustic microscopy and tomography are widely used in laboratory neuroscience, cancer, and drug research to watch blood vessels and test new molecular agents.
\end{itemize}

\section{Advantages}
\begin{itemize}
  \item \textbf{Works Across Multiple Scales:} PAI can zoom in on microscopic surface capillaries or map blood vessels several centimeters deep in soft tissue.
  \item \textbf{Dye-Free Blood and Oxygen Detail:} Natural hemoglobin, melanin, fat, and collagen provide rich vascular and oxygenation maps without needing injected dyes.
  \item \textbf{Safe to Repeat Without Radiation:} Non-ionizing light allows repeated bedside scans to track chronic arthritis, bowel inflammation, or cancer therapy response.
  \item \textbf{Built Right Into Handheld Ultrasound:} Sharing the same ultrasound probe lets doctors see live grayscale anatomy and color blood-oxygen maps together.
  \item \textbf{Flexible Molecular Expansion:} PAI works with approved dyes (ICG and methylene blue) and experimental nanoparticles for lymph-node and targeted molecular studies.
\end{itemize}

\section{Disadvantages}
\begin{itemize}
  \item \textbf{Limited Depth Reach:} Light weakens quickly inside the body, keeping clinical PAI to about 2 to 4~centimeters deep---so it cannot scan deep abdominal organs or the adult brain through the skull.
  \item \textbf{Depth-Related Color Distortion:} Overlying tissues absorb some light colors more than others (spectral coloring), which can make deep oxygen readings less exact.
  \item \textbf{Bone, Air, and Darker Skin Challenges:} Bone and air-filled lungs block sound waves, while higher natural melanin in darker skin absorbs more surface light and can add extra echo clutter.
  \item \textbf{Laser Eye Safety and Higher Cost:} Pulsed lasers require protective eyewear, strict skin-energy limits, ultrasound gel, and more expensive hardware than regular ultrasound.
  \item \textbf{Limited Routine Clinical Availability:} Aside from one FDA-approved adjunctive breast system, most PAI uses are still investigational and need standard calibration, regulatory approval, and large clinical trials.
\end{itemize}

\section{Contraindications}

\subsection{Absolute}
\begin{itemize}
  \item \textbf{No Eye Protection or Unapproved Substitution:} Running the laser without proper protective goggles is strictly forbidden. Outside the FDA-cleared adjunctive breast system, PAI must never replace standard imaging or biopsy.
  \item \textbf{Contrast Dye and Nanoparticle Exclusions:} Indocyanine green must not be used in anyone with a severe allergy to ICG or iodine-containing dyes. Methylene blue must not be used in people with G6PD enzyme deficiency or those taking serotonergic medicines (due to serotonin syndrome risk). Experimental nanoparticles are restricted strictly to safety-vetted FDA/IRB clinical trials.
\end{itemize}

\subsection{Relative}
\begin{itemize}
  \item \textbf{Open Wounds or Dark Tattoos:} Open sores, active skin infections, fresh burns, or dense dark tattoos over the scan spot soak up extra light and create echo clutter, so they are avoided or handled under strict sterile study rules.
  \item \textbf{Deep Targets, Bone, Air, or Very Dark Skin:} Targets deeper than 4~centimeters, areas behind bone or gas-filled bowel, or very high skin melanin can reduce depth and diagnostic clarity.
\end{itemize}

\section{Patient Safety Guidelines}
\begin{itemize}
  \item \textbf{Strict Laser Energy and Skin Protection:} Lasers stay strictly within ANSI~Z136 and IEC~60825 exposure limits, with extra care for darker skin tones and pigmented moles where higher melanin absorbs more heat. Because those exposure limits were established for skin rather than internal organs, investigational endoscopic and catheter-based PAI undergo extra safety testing.
  \item \textbf{Protective Goggles, Room Controls, and Automatic Shutoffs:} High-power Class~IV laser scanners require controlled-access rooms and wavelength-matched laser goggles for both patients and staff (though newer LED-based research prototypes avoid these laser-eye hazards), supported by warning signs and probe sensors that turn off light flashes if the probe lifts off the skin.
  \item \textbf{Ultrasound Safety and Clean Gel:} Paired ultrasound follows AIUM/FDA safety limits (TI and MI), probes are disinfected between patients, and sterile covers and gel are used near open skin, biopsies, or surgery.
  \item \textbf{Safe Dye and Nanoparticle Screening:} Staff screen for dye allergies, G6PD deficiency, and serotonergic drug interactions before giving approved dyes, and restrict experimental nanoparticles strictly to FDA/IRB-regulated trials.
  \item \textbf{Clear Clinical Boundaries:} Routine care is limited to cleared adjunctive breast imaging; all other human PAI scans run under IRB oversight without replacing standard diagnostics.
\end{itemize}

\section{Observing for Abnormal Patterns}
\begin{itemize}
  \item \textbf{Tangled Blood Vessels and Low Oxygen in Tumors:} In adjunctive breast PAI and cancer trials, malignant tumors often show messy internal blood vessels, extra blood around the edges, and lower oxygen levels compared with harmless lumps.
  \item \textbf{Swollen, Low-Oxygen Joint Linings:} In rheumatoid and psoriatic arthritis trials, active joint inflammation shows up as bright blood buildup paired with lower oxygen inside the joint lining---sometimes earlier than standard Doppler ultrasound detects it.
  \item \textbf{Poor Circulation and Fragile Plaques:} In vascular studies, fewer tiny capillaries and lower oxygen signal poor circulation, while bright fat and blood peaks inside neck-artery plaques suggest bleeding-prone buildup.
  \item \textbf{Depth Loss and Color Shifts:} Because light fades beyond 2 to 4~centimeters and surface tissues filter colors unevenly, deep oxygen numbers can drift and falsely look like low oxygen.
  \item \textbf{Angle and Echo Artifacts:} Viewing from one side with a handheld probe can slightly stretch vertical vessels or add streaks, while dark skin pigment, thick hair follicles, or metal biopsy needles can create echoing ``ghost'' signals.
\end{itemize}

\section{Dealing with Uncertainties}
\begin{itemize}
  \item \textbf{Light Fading and Color Distortion:} Rapid light loss past 2 to 4~centimeters lowers image clarity, and overlying blood or fat alters the light colors reaching deeper spots, skewing oxygen calculations.
  \item \textbf{Skin-Tone and Echo Clutter:} Higher skin melanin, dense hair, or metal needles generate strong surface sound waves that bounce and can mask deeper tissue or mimic focal spots.
  \item \textbf{Probe Geometry and Bone/Air Blocks:} Handheld probes can distort shapes slightly, while bone and air-filled lungs or bowel block sound waves completely.
  \item \textbf{Overlap Between Inflammation and Cancer:} Harmless inflammation, active arthritis, benign vascular breast lumps, and early cancers can all show high blood signals, limiting standalone certainty.
  \item \textbf{Confirming with Proven Care:} Teams cut down glitches with degassed gel, multi-angle ultrasound pairing, and calibration checks. Uncertain breast findings follow standard BI-RADS imaging and needle biopsy, while research arthritis, melanoma, or vascular results are checked against Doppler ultrasound, MRI, angiography, or surgical pathology.
\end{itemize}

\section{Augmenting Methods}
\begin{itemize}
  \item \textbf{Live Grayscale and Power Doppler Ultrasound:} Built-in ultrasound provides clear anatomical landmarks, lump shapes, and larger blood-vessel flow right beside PAI's tiny-vessel and oxygen maps.
  \item \textbf{Mammography, MRI, and Angiography:} Mammograms spot tiny calcifications, while MRI and CT/MR angiography see deep joints, deep breast tissue, and full blood-vessel trees beyond 2 to 4~centimeters.
  \item \textbf{Inflammation and Circulation Tests:} Blood markers (such as CRP, ESR, rheumatoid factor, and anti-CCP) and simple ankle-brachial blood-pressure checks add important context in arthritis and vascular trials.
  \item \textbf{Biopsy and Lab Pathology:} Core-needle biopsy and lymph-node pathology remain the definitive gold standards for breast lumps, melanoma margins, and nodal checks.
  \item \textbf{Light-Correction and AI Software:} Smarter computer models that correct for depth, skin tone, and echo clutter are improving oxygen accuracy across patients.
\end{itemize}

\section{Comparison with Alternative Modalities}
\begin{itemize}
  \item \textbf{Versus Standard Ultrasound and Power Doppler:} PAI sees tiny capillary blood and oxygen levels that are too slow for power Doppler to pick up, while standard ultrasound sees deeper, costs less, needs no laser goggles, and is universally approved.
  \item \textbf{Versus Contrast-Enhanced MRI:} MRI sees deep joints, full breasts, and entire blood-vessel networks without depth limits, while handheld PAI offers fast, radiation-free, gadolinium-free bedside checks at lower cost.
  \item \textbf{Versus Pure Optical Microscopy (OCT / Confocal):} PAI sees centimeters deep instead of stopping at 1~millimeter, while pure optical tools offer finer cell-level detail right at the surface.
  \item \textbf{Versus Mammography:} PAI avoids X-rays and breast squeezing and works well in dense breasts, while mammography remains the primary whole-breast screening standard for microcalcifications.
\end{itemize}

\section{Key References}
\begin{itemize}
  \item Wang LV, Hu S. Photoacoustic tomography: in vivo imaging from organelles to organs. \textit{Science}. 2012;335(6075):1458--1462.
  \item Steinberg I, Huland DM, Vermesh O, Frostig HE, Tummers WS, Gambhir SS. Photoacoustic clinical imaging. \textit{Photoacoustics}. 2019;14:77--98.
  \item Attia ABE, Balasundaram G, Moothanchery M, et al. A review of clinical photoacoustic imaging: current and future trends. \textit{Photoacoustics}. 2019;16:100144.
  \item American National Standards Institute (ANSI). \textit{ANSI Z136.1: American National Standard for Safe Use of Lasers}. Orlando, FL: Laser Institute of America, 2022.
  \item Riksen JJM, Nikolaev AV, van Soest G. \textit{Photoacoustic imaging on its way toward clinical utility: a tutorial review focusing on practical application in medicine}. Journal of Biomedical Optics (National Library of Medicine, PMC). Available at: \url{https://pmc.ncbi.nlm.nih.gov/articles/PMC10249868/}.
\end{itemize}

\section{Further Reading}
\begin{itemize}
  \item Wikipedia contributors. \textit{Photoacoustic imaging}. Wikipedia, The Free Encyclopedia. Available at: \url{https://en.wikipedia.org/wiki/Photoacoustic_imaging}.
\end{itemize}
